\chapter{Nociones básicas de análisis complejo}\label{appendix1}

\section{Funciones holomorfas}

\subsection{Derivada compleja}

\transclude{fn-compleja-derivable-pnt}

\transclude{fn-holomorfa}

\transclude{fn-holomorfa-pnt}

\transclude{fn-entera}

\transclude{prop-fn-compleja-derivable}%DEMOSTRACIÓN

\subsection{Ecuaciones de Cauchy-Riemann}

\transclude[teo]{teo-cauchy-riemann}

\transclude[prop]{prop-consecuencias-cauchy-riemann}

\section{Operadores de Wirtinger y el operador laplaciano}

\subsection{Derivadas de Wirtinger}

\transclude{operadores-wirtinger}

También podemos expresar las derivadas parciales habituales en términos de las derivadas de Wirtinger:

\transclude{lem-derivadas-parciales-wirtinger}

\transclude{fn-antiholomorfa}

\transclude{obs-holomorfa-conjugada-antiholomorfa}

\transclude{prop-fn-holomorfa-iff-wirtinger}

\transclude{obs-derivada-holomorfa-wirtinger}

\transclude{prop-regla-cadena-wirtinger}

\transclude{cor-wirtinger-composicion-holomorfas}%DEMOSTRACIÓN

\transclude{cor-wirtinger-composicion-antiholomorfas}%DEMOSTRACIÓN

\subsection{El operador laplaciano}

\begin{defn}[Laplaciano]\label{defn:laplaciano}
    En $\mathbb{D}$, el operador laplaciano $\Delta \colon \mathcal{C}^2 \left(\mathbb{D}\right) \to \mathcal{C}\left(\mathbb{D}\right)$ viene dado por
    \[
    \forall z \in \mathbb{D} : \Delta f(z) = \pdv[2]{f}{x}(z) + \pdv[2]{f}{y}(z) \quad \tex{donde } z = x + iy. 
    \]
\end{defn}

\transclude{lem-laplaciano-wirtinger}

\transclude{lem-laplaciano-cambio-variable}

\section{Transformaciones de Möbius}

\transclude{transformacion-mobius}

\transclude[prop]{prop-transformacion-mobius-composicion}

\transclude{transformacion-mobius-circunferencias-generalizadas}

\section{Teorema de la función inversa y de la aplicación abierta}

\transclude[teo]{teo-fn-inversa-holomorfas}

\transclude{teo-apl-abierta-compleja}



\transclude{inyectividad-local}

\transclude{lem-localmente-inyectiva-implica-inversa-no-nula}%DEMOSTRACIÓN



\section{Integración compleja}

\transclude{curva-topologica}

% \transclude{curva-cerrada}

% \transclude{curva-simple}

% \transclude{curva-jordan}

\transclude{camino}

% \transclude{longitud-camino}

% \transclude{camino-opuesto}

\transclude{integral-linea-compleja}

\transclude{integral-linea-compleja-longitud}



\transclude[teo]{teo-cauchy-goursat}

\transclude[teo]{teo-formula-integral-cauchy-disco}

\begin{cor}[Propiedad del valor medio para funciones holomorfas]\label{teo:valor-medio-fn-holomorfa}\mbox{}\\
    Sea $\Omega \subset \C$ abierto, $f \in \mathcal{H}(\Omega)$, $z_0 \in \Omega$ y $r > 0$ tales que $\overline{D(z_0,r)} \subset \Omega$. Entonces,
    \[
    f(z_0) = \frac{1}{2\pi} \int_0^{2\pi} f\left(z_0 + re^{it}\right) \, \odif{t}.
    \]
\end{cor}
\begin{dem}
    Por la \exhyperref[teo:formula-integral-cauchy-disco]{teo-formula-integral-cauchy-disco}{fórmula integral de Cauchy para discos},
    \[
    f(z_0) = \frac{1}{2\pi i} \int_{\abs{\zeta-z_0}=r} \frac{f(\zeta)}{\zeta-z_0} \, \odif{\zeta}.
    \]
    Parametrizando $\zeta(t) = z_0 + re^{it}$, $t \in [0,2\pi]$, se tiene $\odif{\zeta} = ire^{it} \, \odif{t}$ y $\zeta-z_0 = re^{it}$:
    \[
    f(z_0) = \frac{1}{2\pi i} \int_0^{2\pi} \frac{f\left(z_0+re^{it}\right)}{re^{it}} \, ire^{it} \, \odif{t} = \frac{1}{2\pi} \int_0^{2\pi} f\left(z_0+re^{it}\right) \, \odif{t}.
    \]

    \vspace{-\baselineskip}
    \vspace{-\belowdisplayskip}
\end{dem}

\transclude[teo]{teo-formula-integral-cauchy-derivadas}




\transclude[teo]{teo-ceros-aislados}
\vspace{0.3em}

\begin{cor}\label{cor:fn-holomorfa-compacto-imp-finitud-ceros}
    Sea $\Omega \subset \C$ un \exhyperref[defn:dominio]{dominio}{dominio}, $K \subset \Omega$ un \exhyperref[defn:compacidad]{compacidad}{compacto} y sea $f \in \exhyperref[defn:fn-holomorfa]{fn-holomorfa}{\mathcal{H}}(\Omega)$ no identicamente nula. Entonces, $f$ tiene un número finito de ceros en $K$.
\end{cor}
\begin{dem}
    Supongamos por contradicción que existe una sucesión $\left(z_n\right)_{n \in \N} \subset K$ tal que $\forall n \neq m : z_n \neq z_m \we f(z_n) = 0$. Como $K$ es compacto, $\exists \left(z_{n_k}\right)_{k \in \N} \subset \left(z_n\right)_{n \in \N}$ tal que $z_{n_k} \xrightarrow{k \to \infty} z_0 \in K$. Entonces, por el \exhyperref[teo:ceros-aislados]{teo-ceros-aislados}{principio de los ceros aislados}, $f \equiv 0$ en $\Omega$. Contradicción.
\end{dem}

\begin{cor}[Principio de identidad]\label{cor:principio-identidad}
    Sea $\Omega \subset \C$ un \exhyperref[defn:dominio]{dominio}{dominio}, sean $f, g \in \mathcal{H}(\Omega)$ y $\left(a_n\right)_{n\in\N} \subset \Omega$ tales que $a_n \xrightarrow{n\to\infty} a \in \Omega \we \forall n \in \N : f(a_n) = g(a_n)$
    \[
    \implies f \equiv g \tex{ en } \Omega.
    \]
\end{cor}

\transclude[teo]{teo-modulo-maximo}

\transclude[cor]{cor-modulo-maximo}

\transclude{prop-modulo-constante-imp-constante}

\transclude{orden-cero-fn-holomorfa}

\transclude{valencia-fn-holomorfa}

\transclude{singularidad-aislada}

\transclude{singularidad-evitable}

\transclude{polo}

\transclude{singularidad-esencial}

\transclude{teo-singularidad-evitable-riemann}

\transclude{fn-meromorfa}

\chapter{Ejercicios}

\section{Ejercicios del capítulo 1}

\begin{ejer}\label{ejer:identidad-auxiliar-disco-unidad}
    Sean $\alpha, \beta \in \mathbb{D}$. Demuestra la siguiente identidad:
    \[
    1 - \abs{\frac{\alpha - \beta}{1 - \overline{\beta} \alpha}}^2 = \frac{(1 - \abs{\alpha}^2)(1 - \abs{\beta}^2)}{\abs{1 - \overline{\beta} \alpha}^2}.
    \]
\end{ejer}
\begin{sol}
    Partiendo del lado izquierdo, tenemos que
    \[
    1 - \abs{\frac{\alpha - \beta}{1 - \bar{\beta} \alpha}}^2 = \frac{\abs{1 - \bar{\beta} \alpha}^2 - \abs{\alpha - \beta}^2}{\abs{1 - \bar{\beta} \alpha}^2}.
    \]
    En el numerador, desarrollando ambos términos, obtenemos
    \[
    \begin{aligned}
        \abs{1 - \bar{\beta} \alpha}^2 - \abs{\alpha - \beta}^2 &= \left(1 - \bar{\beta} \alpha\right)\left(1 - \beta \bar{\alpha}\right) - (\alpha - \beta)(\bar{\alpha} - \bar{\beta}) \\
        &= 1 - \cancel{\bar{\alpha} \beta} - \cancel{\alpha \bar{\beta}} + \abs{\alpha}^2 \abs{\beta}^2 - \left(\abs{\alpha}^2 - \cancel{\alpha \bar{\beta}} - \cancel{\bar{\alpha} \beta} + \abs{\beta}^2\right) \\
        &= 1 - \abs{\alpha}^2 - \abs{\beta}^2 + \abs{\alpha}^2 \abs{\beta}^2 = (1 - \abs{\alpha}^2)(1 - \abs{\beta}^2).
    \end{aligned}
    \]
    Por tanto, se sigue la identidad deseada.\hfill \qedsymbol
\end{sol}

\begin{ejer}\label{ejer:desigualdad-schwarz-pick}
    Sea $f \in \exhyperref[defn:clase-schur]{clase-schur}{\mathcal{S}}$. Demuestra que
    \[
    \forall z, w \in \mathbb{D} : \abs{\frac{f(z) - f(w)}{z-w}}^2 \leq \frac{1 - \abs{f(z)}^2}{1 - \abs{z}^2} \cdot \frac{1 - \abs{f(w)}^2}{1 - \abs{w}^2}.
    \]
\end{ejer}
\begin{sol}
    Como $f \in \mathcal{S}$, por el \exhyperref[teo:schwarz-pick]{teo-schwarz-pick}{teorema de Schwarz-Pick}, tenemos que
    \[
    \abs{\frac{f(z) - f(w)}{1 - \overline{f(w)} f(z)}} \leq \abs{\frac{z - w}{1 - \bar{w} z}} \implies 1 - \abs{\frac{z - w}{1 - \bar{w} z}}^2 \leq 1 - \abs{\frac{f(z) - f(w)}{1 - \overline{f(w)} f(z)}}^2. 
    \]

    Por el \excref[ejer:identidad-auxiliar-disco-unidad]{ejer-identidad-auxiliar-disco-unidad}, esta desigualdad se traduce a
    \[
    \begin{aligned}
        \frac{(1 - \abs{z}^2)(1 - \abs{w}^2)}{\abs{1 - \bar{w} z}^2} &\leq \frac{(1 - \abs{f(z)}^2)(1 - \abs{f(w)}^2)}{\abs{1 - \overline{f(w)} f(z)}^2} \\
        \implies \abs{\frac{1 - \bar{f(w)} f(z)}{1 - \bar{w} z}}^2 &\leq \frac{1 - \abs{f(z)}^2}{1 - \abs{z}^2} \cdot \frac{1 - \abs{f(w)}^2}{1 - \abs{w}^2}.
    \end{aligned}
    \]

    Finalmente, observamos que:
    \[
    \begin{aligned}
        \abs{\frac{f(z) - f(w)}{z - w}}^2 &= \abs{\frac{f(z) - f(w)}{1 - \overline{f(w)} f(z)}}^2 \cdot \abs{\frac{1 - \overline{f(w)} f(z)}{1 - \bar{w} z}}^2 \cdot \abs{\frac{1 - \bar{w} z}{z - w}}^2 \\
        &\leq 1 \cdot \frac{1 - \abs{f(z)}^2}{1 - \abs{z}^2} \cdot \frac{1 - \abs{f(w)}^2}{1 - \abs{w}^2} \cdot 1 = \frac{1 - \abs{f(z)}^2}{1 - \abs{z}^2} \cdot \frac{1 - \abs{f(w)}^2}{1 - \abs{w}^2}.
    \end{aligned}
    \]

    \vspace{-\baselineskip}
    \vspace{-\belowdisplayskip}
    \hfill \qedsymbol
\end{sol}

\section{Ejercicios del capítulo 2}

\begin{ejer}
    Demuestra que la métrica de Poincaré $\wp$ satisface que
    \[
    \forall z, w \in \mathbb{D} : \wp(z, w) = 2 \tanh^{-1} \abs{\frac{z - w}{1 - \overline{w}z}}.
    \]
\end{ejer}
\begin{sol}
    Recordemos la definición de la tangente hiperbólica:
    \[
    \tanh(z) = \frac{\sinh(x)}{\cosh(x)} = \frac{e^x - e^{-x}}{e^x + e^{-x}} = \frac{e^{2x} - 1}{e^{2x} + 1}.
    \]
    Entonces, despejando, tenemos que $e^{2x} \left(1 - \tanh(x)\right) = 1 + \tanh(x)$, luego
    \[
    e^{2x} = \frac{1 + \tanh(x)}{1 - \tanh(x)} \implies \tanh^{-1}(y) = \frac{1}{2} \, \log \frac{1 + y}{1 - y}.
    \]
    En particular, para $y = \varrho (z, w)$ con $z, w \in \mathbb{D}$, tenemos que
    \[
    \tanh^{-1} \left(\varrho(z, w)\right) = \frac{1}{2} \, \log \frac{1 + \varrho(z, w)}{1 - \varrho(z, w)} = \frac{1}{2} \, \wp(z, w)
    \]
    por definición de la métrica de Poincaré. Por tanto, como $\varrho(z, w) = \abs{\frac{z - w}{1 - \overline{w} z}}$, se tiene que
    \[
    \implies \wp(z, w) = 2 \, \tanh^{-1} \left(\varrho(z, w)\right) = 2 \, \tanh^{-1} \abs{\frac{z - w}{1 - \overline{w} z}}.
    \]

    \vspace{-\baselineskip}
    \vspace{-\belowdisplayskip}
    \hfill\qedsymbol
\end{sol}

\begin{ejer}
    Sea $f \colon \Omega_1 \to \Omega_2$ una aplicación holomorfa entre dos dominios $\Omega_1, \Omega_2 \subset \C$ dotados de densidades métricas $\mu_1, \mu_2$ respectivamente. Demuestra que
    \[
    f \tex{ es una isometría entre } (\Omega_1, d_{\mu_1}) \tex{ y } (\Omega_2, d_{\mu_2}) \iff \forall z \in \Omega_1 : (f^* \mu_2)(z) = \mu_1(z).
    \]
\end{ejer}
% \begin{sol}
%     Veamos ambas implicaciones por separado:
%     \begin{itemize}
%         \item[$\boxed{\Rightarrow}$]
%         \item[$\boxed{\Leftarrow}$] 
%     \end{itemize}
% \end{sol}
%DEMOSTRACIÓN

\begin{ejer}%EJERCICIO 
    Demuestra que la curvatura de la métrica esférica es constante e igual a $+1$.
\end{ejer}
% \begin{sol}

% \end{sol}

%COMPLETAR mirar Geometry: Euclid and Beyond.


\chapter{Convergencia y continuidad uniformes}

\begin{prop}\label{prop:convergencia-uniforme-continuidad-uniforme}
    Sean $X, Y$ \exhyperref[defn:esp-metrico]{esp-metrico}{espacios métricos}, $f \colon X \to Y$ y sea $\left(f_n\right)_{n \in \N}$ una sucesión de funciones $f_n \colon X \to Y$ \exhyperref[defn:continuidad-uniforme]{continuidad-uniforme}{uniformemente continuas} que \exhyperref[defn:convergencia-uniforme]{convergencia-uniforme}{converge uniformemente} a $f$. Entonces, $f$ es uniformemente continua.
\end{prop}
\begin{dem}
    Sea $\varepsilon > 0$. Entonces, $\exists N \in \mathbb{N} : \forall x \in X : d_Y(f_N(x), f(x)) < \frac{\varepsilon}{3}$.

    Por otro lado, como $f_N$ es uniformemente continua, $\exists \delta > 0 : \forall x_1, x_2 \in X : d_X(x_1, x_2) < \delta \implies d_Y(f_N(x_1), f_N(x_2)) < \frac{\varepsilon}{3}$. Por tanto, si $d_X(x_1, x_2) < \delta$, por la \exhyperref[defn:metrica:desigualdad-triangular]{metrica}{desigualdad triangular}, se tiene que
    \begin{align*}
        d_Y(f(x_1), f(x_2)) &\leq d_Y(f(x_1), f_N(x_1)) + d_Y(f_N(x_1), f_N(x_2)) + d_Y(f_N(x_2), f(x_2)) \\
        &< \frac{\varepsilon}{3} + \frac{\varepsilon}{3} + \frac{\varepsilon}{3} = \varepsilon.
    \end{align*}
    Luego $f$ es uniformemente continua.
\end{dem}